% Source: https://artofproblemsolving.com/community/c3h1223478_2004_mathcounts_nationals
% Stable identity: aops-topic-1223478-q02
$C^E=EC$ is a two-digit number ending in $C$ whose tens digit is the exponent $E$: trying small cases, $5^2=25$ works ($C=5$, $E=2$), and no other pair does ($E\ge3$ makes $C^E$ too large or of the wrong form; e.g. $6^2=36$ has tens digit $3\ne2$). So $E=2$, $C=5$.

$AD^2=AHJ$ is a three-digit square whose first digit equals the first digit of its root: $AD\in\{10,\dots,31\}$ with $A\in\{1,2,3\}$. $B^D=HB$ is a two-digit power ending in its base: candidates $4^3=64$, $6^2=36$, $6^3=216$ (too big), $9^3=729$ (too big), $5^3$ excluded ($C=5$). $D=2$ is taken by $E$, so $B=4$, $D=3$, $H=6$. Then $AD^2=A3^2$ must have hundreds digit $A$ and tens digit $6$: $13^2=169$ works ($A=1$, $J=9$); $23^2=529$ and $33$ are not allowed. Finally $AB\cdot F=14F=JG=9G$: $14\cdot7=98$, so $F=7$, $G=8$. All nine digits $1,4,5,3,2,7,8,6,9$ are distinct.

Hence $G^E=8^2=\boxed{64}$.
